\subsection{Three neutral components of conformal extensions}
\label{sec:conformal-models}
The same construction realizes three further census rings. We write
$\mathcal B(\mathfrak g,k)$ for the principal unitary WZW category and
use finite Dynkin labels for its simple objects.

\begin{theorem}\label{thm:conformal-models}
The rings $\RR{4}{22}$, $\RR{4}{91}$ and $\RR{4}{629}$ admit unitary
categorifications. They are the neutral components of the module
categories for the conformal embeddings
\[
 (B_2)_7\subset(D_7)_1,\qquad
 (C_3)_4\subset(B_{10})_1,\qquad
 (B_2)_{12}\subset(E_8)_1,
\]
respectively.
\end{theorem}

\paragraph{Common construction and identification principle.}
A unitary conformal extension supplies an actual connected commutative
separable algebra $A$ in its source category $\mathcal B$.
Consequently $\mathcal E=\mathcal B_A$ is a unitary fusion category,
$F:\mathcal B\to\mathcal E$, $U\mapsto U\otimes A$, is a central
 tensor functor, and the local modules are the level-one target category
$\mathcal L$. These facts use the conformal algebra, not just a modular
invariant \cite[Sections~6.1--6.2]{DMNO}.
In each case $A$ belongs to the even part of the source parity grading.
The two homogeneous parts of any module are submodules. Thus simples
have definite parity, relative tensor product respects parity, and odd
free modules make the grading faithful. Its neutral component
$\mathcal C=\mathcal E_0$ is therefore an actual unitary fusion category.

We identify its full based ring from free induction and local modules.
Two standard tests resolve the few possible tensor completions. First,
\cite[Corollary~3.30]{DMNO} and \cite[Proposition~8.23.11]{EGNO} give
\begin{equation}\label{eq:conformal-neutral-center}
 Z(\mathcal E)\simeq\mathcal B\boxtimes\mathcal L^{\rm rev},
 \qquad
 Z(\mathcal C)\simeq
 \bigl(\Rep(C_2)'\subset Z(\mathcal E)\bigr)_{\Rep(C_2)},
\end{equation}
where $\Rep(C_2)$ is the canonical grading subcategory. If a commutative
candidate ring has formal codegree $f$, the actual known center must
contain a simple of dimension $\dim(\mathcal C)/f$
\cite[Theorem~2.13]{Ostrik3}. Second, indicators can be computed directly
from the actual center and its forgetful functor \cite{NS}.
These tests concern the constructed category and do not presume that
archived modular data have been categorified.

\paragraph{The $B_2$ level-seven component.}
The source \cite[Section~3.5, Figure~11(a)]{CRT} gives
\[
 A=(0,0)\oplus(1,6)\oplus(2,2)\oplus(5,0),
 \qquad \mathcal L=\mathcal B(D_7,1).
\]
Parity is the second Dynkin coordinate modulo two. The six even vertices,
in source order $(1,9,10,2,5,6)$, will be denoted $(1,x,y,g,e,f)$.
Here $e=F(1,0)$ and $g=F(7,0)$ are simple, hence central in the fusion
ring, and $e$ is self-dual. The published vector-action graph gives
\[
 E=N_e=\begin{pmatrix}
 0&0&0&0&1&0\\
 0&1&1&0&0&1\\
 0&1&1&0&0&1\\
 0&0&0&0&1&0\\
 1&0&0&1&0&2\\
 0&1&1&0&2&2
 \end{pmatrix}.
\]
The local category has fusion group $C_4$; an odd local spinor $z$
satisfies $z^2=g$ \cite[Theorem~7.1]{Bock}. The odd simple
$s=F(0,1)$ is self-dual and central, and $zs$ is one of $x,y$.
Consequently $(zs)^*=gzs$: duality on this pair is the action of $g$.
Dimensions force $g$ to fix $e,f$ and either fix or interchange $x,y$.
The first row of $E^2$ gives $e^2=1+g+2f$; the matrix
$(E^2-I-N_g)/2$ is integral only for the interchange. Thus $x^*=y$.

Put $G=N_g$, $H=(E^2-I-G)/2$ and $X=N_x$.
The equation $ef=x+y+2e+2f$ and rigidity imply
\[
 X+GX=EH-2E-2H,\qquad X^{\mathsf T}=GX.
\]
Also $X$ commutes with $E,G,H$, and its unit row is the basis vector
of $x$. These are linear equations over $\QQ$ in the $36$ entries of
$X$. Their coefficient matrix has rank $36$. The unique solution is
nonnegative integral and, with $N_y=GX$ and $N_f=H$, gives exactly all
$216$ coefficients of $\RR{4}{22}$. The accompanying checker extracts
the graph directly from the source PDF and reconstructs this linear
system, its rank, and the full tensor comparison.

\paragraph{The $B_2$ level-twelve component.}
Here \cite[Section~3.6, Figure~13]{CRT} gives
\begin{align*}
 A={}&(0,0)\oplus(0,6)\oplus(1,8)\oplus(3,2)\oplus(3,8)\\
 &\oplus2(4,4)\oplus(6,6)\oplus(7,2)\oplus(12,0).
\end{align*}
The target category $\mathcal L$ is $\Vec$, so $A$ is Lagrangian and
$Z(\mathcal E)\simeq\mathcal B(B_2,12)$.
The even graph component has source vertices $(1,3,7,8,4,6)$.
The central simple $x=F(1,0)$ has multiplication matrix of characteristic
polynomial $t^2(t^4-5t^3+10t-5)$. Every eigenvalue has multiplicity at
most two. A noncommutative matrix summand $M_n(\CC)$ of the complexified
fusion ring would contribute at least $n^2\ge4$ to the multiplicity of
an eigenvalue of this central element. The ring is therefore commutative.

The only duality permutations fixing $1,x$ and commuting with $N_x$
are the identity and the interchange of vertices $7,8$. In the latter
case, Frobenius reciprocity and associativity with $x$ force coefficients
$1/2$ and $3/2$, which are impossible. For identity duality the same
linear constraints leave one parameter $t$. In the indicated basis
$(1,x,u,v,w,z)$,
\[
 N_u=\begin{pmatrix}
 0&0&1&0&0&0\\
 0&0&1&1&1&1\\
 1&1&4-t&t-1&1&2\\
 0&1&t-1&3-t&2&2\\
 0&1&1&2&2&2\\
 0&1&2&2&2&3
 \end{pmatrix}.
\]
All other products are determined and associative; integrality and
nonnegativity give $t\in\{1,2,3\}$. The middle value is exactly
$\RR{4}{629}$. The outer values have formal codegrees $3,15$.

The unique nontrivial invertible $J=(12,0)$ in $\mathcal B(B_2,12)$
is a boson. Thus \eqref{eq:conformal-neutral-center} identifies the
actual center as $(\mathcal B(B_2,12)_{\rm ad})_{\langle J\rangle}$.
Among the $49$ even weights, the involution
$(a,b)\mapsto(12-a-b,b)$ has $21$ ordinary orbits and seven fixed
points. Each ordinary orbit gives one simple of the original Weyl
dimension; each fixed point gives two simples of half that dimension.
This determines all $35$ simple dimensions of the center. With
$D=\dim(A)/2$, exact Weyl arithmetic in $\QQ(\zeta_{60})$ shows that
neither $D/3$ nor $D/15$ occurs. The formal-codegree test excludes both
outer completions, proving the required tensor identification.

\paragraph{The $C_3$ level-four component.}
The adjoint conformal branching theorem
\cite[Theorem~3.9 and Proposition~4.1]{CKMP} gives the three local
Ising modules $1,\psi,\sigma$ with underlying source objects
\begin{align*}
 U(1)=A&=000\oplus210\oplus004\oplus121,\\
 U(\psi)&=200\oplus030\oplus022\oplus301,\\
 U(\sigma)&=2(111).
\end{align*}
Their dimensions are $1,1,\sqrt2$ and their twists are
$1,-1,\zeta_{16}^{5}$. The source parity is $a+c$ modulo two;
there are $19$ even weights.
For free modules the exact induction Gram matrix is
\[
 G_{\lambda\mu}=\langle F(\lambda),F(\mu)\rangle
      =\sum_\alpha[A:\alpha]N_{\lambda\alpha}^{\mu}.
\]
The three local columns are known from branching, while
$\tau=F(010)$ has norm one and is orthogonal to them. Subtracting these
four columns leaves $2vv^{\mathsf T}$ with $v\ge0$ integral and
$v_{002}=1$. The row of norm two must be a sum of two distinct unit
vectors. Rank one forces every other residual row to be proportional
to it. Thus there are exactly two further simples $x,y$, with identical
restriction columns, and $F(002)=x+y$.
Dominance of free induction rules out further simples.

The free image has rank five in the six-dimensional complexified fusion
ring. It is central, so the ring is commutative: a noncommutative matrix
block would make the difference between algebra and center dimensions
at least three. The objects $1,\psi,\sigma,\tau$ are self-dual, leaving
only the alternatives $x^*=y$ and self-duality of both $x,y$.
Source fusion, the induction matrix, the Ising products, Frobenius
reciprocity and associativity give a bounded exact completion problem.
Its four or five free integer parameters lie between zero and four,
as follows from coefficient nonnegativity. In the non-self-dual case
there is one tensor, equal in all $216$ entries to $\RR{4}{91}$.
There are three self-dual possibilities, with respectively
\[
 x^2=y^2=1+2y+2\tau,\qquad
 1+x+y+2\tau,\qquad
 1+2x+2\tau.
\]

To exclude the middle possibility, note that a self-dual simple $X$
has a symmetric self-duality on $X\otimes X$. Indeed, if the pivotal
transpose of $f:X\to X^*$ is $\epsilon f$, then the transpose of
$f\otimes f$, under the reversed dual-tensor identification, is
$\epsilon^2(f\otimes f)$. Every self-dual constituent of multiplicity
one therefore has second indicator $+1$. The middle rule would give
$\nu_2(x)=\nu_2(y)=1$.
On the other hand, the known center
$Z(\mathcal E)\simeq\mathcal B(C_3,4)\boxtimes\mathcal L^{\rm rev}$
and the displayed branching give
\[
 \nu_2(F(\mu))=
 \frac1{\dim\mathcal E}\sum_{\lambda,a,\nu}
 \theta_\lambda^2\theta_a^{-2}d_\lambda d_a
 [U(a):\nu]N_{\nu\lambda}^{\mu}.
\]
Exact evaluation in $\QQ(\zeta_{32})$ gives
$\nu_2(F(002))=0$, contradicting $\nu_2(x+y)=2$.

For the outer two completions the formal codegrees include
$8\pm4\sqrt2$. The boson $J=004$ belongs to $A$, so $F(J)=1$;
the canonical grading subcategory in the known center is generated by
$(J,1)$. Formula~\eqref{eq:conformal-neutral-center} gives
\[
 Z(\mathcal C)\simeq
 (\mathcal B(C_3,4)_{\rm ad})_{\langle J\rangle}
 \boxtimes\mathcal L^{\rm rev}.
\]
The action $J(a,b,c)=(b,a,4-a-b-c)$ has eight ordinary orbits and three
fixed weights among the even weights. Fixed-point splitting gives
rank $14$ for the first factor and rank $42$ for the center.
Its dimension list contains neither
$\dim(\mathcal C)/(8+4\sqrt2)$ nor
$\dim(\mathcal C)/(8-4\sqrt2)$, as checked exactly by the Weyl formula.
The formal-codegree test excludes both outer completions. Only the
non-self-dual target remains, completing the proof of
Theorem~\ref{thm:conformal-models}.

The supplement reconstructs the branching from the abelian ideals of
$C_3$, certifies the source fusion coefficients by exact Weyl-character
identities, and checks every finite completion and dimension comparison.
All three arguments establish categorical existence and full tensor
identifications; they do not classify associators on these rings.
